\documentclass[10pt]{article}
\usepackage[margin=20mm]{geometry}
\usepackage{amsmath,amssymb,fontspec,unicode-math,xurl,hyperref}
\IfFontExistsTF{DejaVu Serif}{\setmainfont{DejaVu Serif}}{\setmainfont{Latin Modern Roman}}
\setmathfont{latinmodern-math.otf}
\setlength{\parindent}{0pt}\setlength{\parskip}{6pt}\sloppy
\begin{document}
Scientific Verification Report: 0085\par
Dynamic Mixture of Curriculum LoRA Experts for Continual Multimodal Instruction Tuning\par
I25-43454ICML 2025, PMLR 267:19011-19033; publisher-linked 23-page PDF\par
Inspection outcome: Major: the norm theorem and the published equal-parameter rationale fail\par
\section*{1. Identity, outcome and scope}
Authors: Chendi Ge, Xin Wang, Zeyang Zhang, Hong Chen, Jiapei Fan, Longtao Huang, Hui Xue and Wenwu Zhu. The controlling 23-page PDF has SHA-256 `3a800ba288be998f7165ac570c5ea98fe633fd54f2519cea158dd89861a7f5ef`; its original bytes were preserved. Page references below are physical PDF pages.\par
The first confirmed Major finding concerns Theorem 3.1 / Theorem K.1. Its assumptions permit equal gradient norms at every layer even though at least one layer has non-collinear expected gradients, so its asserted existence of a layer with different gradient norms is false. A complete token cross-entropy example, compatible transformer realization, and independent arithmetic establish the failure. This damages the claimed formal justification for task-dependent magnitude differences and allocation, without establishing that the benchmark scores or practical heuristic are false.\par
A second Major finding concerns the published equal-parameter rationale: its rank-times-location calculation is not invariant to unequal matrix costs, and the task-1 allocation in Figure 6 contradicts the asserted equality under the paper’s stated targets. This invalidates the stated justification for interpreting the central comparison as matched in trainable parameters. Exact historical checkpoint counts and any effect on Table 2 remain Needs Evidence. A separate release-level routing mismatch also remains distinct from an established error in the published scores.\par
Coverage includes the main method, Theorem 3.1 and Appendix K, evaluation definitions and headline arithmetic, appendices on implementation, routing, thresholds, runtime and budgets, the published allocation heatmap, and static inspection of the official code. No author code was executed or imported, and no training, inference or benchmark reproduction was performed.\par
\section*{2. The formal claim and its actual assumptions}
Theorem 3.1 (p. 3) and Theorem K.1 (p. 22) assume that two tasks have distinct input-output distributions and require different optimal parameter configurations (K.2). They also assume that at least one layer has expected gradients that cannot be related by any real scalar multiplier (K.3). The conclusion is that some layer has different norms of these expected gradients.\par
\[\exists l^*: \left\|\mathbb E_{T_A}[\nabla_{W_{l^*}}\mathcal L]\right\|_2\ne\left\|\mathbb E_{T_B}[\nabla_{W_{l^*}}\mathcal L]\right\|_2.\] The printed formulas and Definition K.4 use the norm of an expected gradient, not the expectation of individual gradient norms. The decisive pages were inspected as images to avoid a formula-extraction error. Both task expectations are evaluated at a common parameter checkpoint, consistent with the shared pretrained-model proxy in Eq. (3). The statement supplies no convergence condition, optimizer constraint, or temporal quantifier that would force a norm inequality.\par
\section*{3. Complete counterexample with finite optima}
Use one fixed image-and-text instruction input and three possible next tokens. Task A has probabilities q\_A=(1/2,1/3,1/6); task B has q\_B=(1/3,1/2,1/6). These can be realized as six-example datasets with token counts (3,2,1) and (2,3,1). Both distributions have full support and are distinct. With logits w in R\textasciicircum{}3, the ordinary expected token cross-entropy is:\par
\[F_q(w)=\log\!\left(\sum_{j=1}^3e^{w_j}\right)-q^\top w,\qquad \nabla F_q(w)=\operatorname{softmax}(w)-q.\] Each risk has finite minimizers w=log(q)+c*1. The optimal predictive distribution is uniquely q. Because q\_A differs from q\_B, no single parameter configuration can achieve both predictive optima on the same input. K.2 is therefore satisfied; uniqueness of every redundant neural parameter is not required.\par
\[g_A=\left(-\frac16,0,\frac16\right),\quad g_B=\left(0,-\frac16,\frac16\right),\quad \|g_A\|_2^2=\|g_B\|_2^2=\frac1{18}.\] No scalar multiple of g\_B equals g\_A: their first components already rule this out. Thus K.3 holds. Nevertheless, the two norms both equal sqrt(2)/6. This is a counterexample to the asserted implication, rather than merely a missing justification for it.\par
To realize the example within the claimed model class, use a bias-free residual transformer with one active feed-forward output projection P. On the fixed instruction, choose its hidden activation u=e1 with unit norm, and a fixed affine readout that subtracts the residual base and maps the active coordinates to logits w=P u. At P=0, the expected gradient of P is g\_q u\textasciicircum{}T; its flattened L2 norm is exactly ||g\_q||. Earlier feed-forward gradients vanish because their path contains P=0.\par
Other attention branches can have both value and output factors zero. Other feed-forward branches can have both factors zero with an activation satisfying phi(0)=0; use bias-free branches or fixed biases. They are identity residual blocks with zero parameter gradients at this checkpoint. A fixed zero visual-to-predictive projector makes the visual transformer’s gradients zero. The theorem contains no requirement that every modality have nonzero influence for every task. Extra inactive dimensions and blocks can be added without changing the example.\par
Hence the active layer has equal nonzero gradient norms and non-collinear expected gradients, while every other transformer layer has equal zero norms. There is no layer with the claimed inequality. Setting P u=log(q\_A) or log(q\_B) realizes the respective finite optima while leaving the other blocks unchanged. The witness is a formal admissible model, not a claim about the actual InternVL2 weights used in the experiments.\par
Sequential presentation does not exclude this checkpoint. With fixed features and only P trained, choose 0<eta<1 and let w\_pre solve w\_pre=eta(softmax(w\_pre)-q\_A). The map is a contraction because the softmax Jacobian has operator norm at most 1. A finite fixed point exists, and one legitimate task-A gradient step maps w\_pre exactly to zero, where task B may begin. The source imposes no requirement that task A be fully optimized first.\par
\section*{4. Proof failure, impact and valid repair limits}
The proof on p. 23 first places equal-norm, non-collinear gradients on the same sphere in Eq. (26). That geometry is consistent with all preceding assumptions. It then asserts that equal norms would require reconciling incompatible updates and would violate sequential training. No mathematical condition supplies this purported contradiction.\par
\[g_A^\top g_B=\frac1{36}>0,\qquad \cos(g_A,g_B)=\frac12,\qquad D_{-g_A}F_B(0)=-\frac1{36}<0.\] The example also shows that non-collinearity does not imply opposing updates. A sufficiently small descent step for either task reduces both losses locally. Independent reviewer arithmetic gives the following values for a task-A step with eta=0.1:\par
Risk | Before step | After task-A step\par
Task A | 1.0986123 | 1.0931493\par
Task B | 1.0986123 | 1.0959271\par
The supported impact is the failure of the stated formal foundation for task-specific gradient-magnitude discrepancies. The paper prominently uses this theorem in Section 3.1 and its contribution list to motivate nonuniform allocation. Appendix location or an easy wording edit does not reduce that demonstrated mathematical impact. The empirical sensitivity heatmap and successful benchmark comparisons can still be observations even though this universal theorem is false.\par
A valid weaker conclusion is that some expected gradient vectors differ; that follows directly from K.3. It does not prove different magnitudes, different layer rankings, or optimal parameter allocation. Even a common positive rescaling of every layer’s gradient norm between tasks preserves all rankings. A defensible revision would remove or weaken the universal norm theorem and describe the measured sensitivity patterns as empirical motivation. A retained norm-separation theorem requires substantive extra assumptions and a new proof; simply relabeling non-collinearity as conflict is insufficient.\par
\section*{5. Major: the published equal-parameter rationale is invalid}
Appendix A (p. 15) states equal trainable parameter counts, justified by rank 4 on all baseline layers versus rank 8 on half the method’s allocation locations. Eqs. (5), (14) and (15) actually constrain counts of experts/layers. A LoRA adapter on a d\_out by d\_in matrix has r(d\_in+d\_out) trainable weights, so selecting half the locations does not in general select half the parameter cost.\par
Target weight | LoRA parameter cost at rank r\par
LLM fused qkv: 2048 to 4096 | 6144 r\par
LLM attention output: 2048 to 2048 | 4096 r\par
Vision qkv: 1024 to 3072 | 4096 r\par
Vision attention output: 1024 to 1024 | 2048 r\par
These dimensions follow the stated InternVL2-2B model and its pinned official configuration: the LLM has 16 query heads and 8 key/value heads, and the vision attention uses ordinary qkv projection. With the paper’s attention-only target description, rank 4 over all 24 LLM and 24 vision blocks gives 1,572,864 LoRA weights.\par
The published Figure 6 was visually checked and independently sampled at its exact two fill colors. Every task row contains 48 selected attention/output subslots. For task 1, the LLM has 21 first-subslot and 15 second-subslot selections; vision has 11 and 1. Regardless of which subslot is assigned to qkv versus output within each modality, rank 8 yields between 1,638,400 and 1,900,544 weights. This is 4.17\% to 20.83\% above the stated attention-only baseline. The counts are reproducible from the archived source figure; they are not measured checkpoint tensors.\par
Thus the published count-based explanation does not establish the claimed equality, and its illustrated task-1 allocation is inconsistent with that equality under the stated target matrices. The public allocator likewise budgets by score-row count and enforces a 25\% minimum vision share rather than weighting actual parameter costs. However, the released baseline wrapper also includes feed-forward target modules when no architecture filter is supplied, unlike the manuscript’s attention-only description. This prevents treating the preceding theoretical counts as verified counts of the original experimental runs.\par
The demonstrated failure of this equality rationale is Major because Appendix A uses it to justify the fairness of the central parameter-constrained comparison. This conclusion does not require an empirical ranking reversal or a claim that the released source produced the original experiments. Actual matched-run counts remain Needs Evidence: supply per-task architecture files and tensor counts for every baseline, distinguishing new trainable adapters, frozen accumulated adapters and auxiliary autoencoders. A cost-constrained allocation can use sum\_j I\_j r\_j(d\_in,j+d\_out,j)<=B. If that changes the selected architecture, its performance requires a new comparison. Table 11 already reports O-LoRA at nominal 2x and 4x budgets still below D-MoLE; that preserves relevant performance counterevidence, but does not make the published equality calculation valid.\par
\section*{6. Official release: routing differs from the manuscript}
The official gcd19/D-MoLE repository was pinned at commit `525fe2b788da85ebe3f5f102c4a77b3d37348c18`. Relevant source bytes and Git blob identities were preserved. Eq. (10) defines a threshold-passing candidate set separately for each sample, and Eq. (12) selects up to two lowest-error eligible task experts for that sample. The training description selects a prior expert together with the current task’s new expert.\par
The public get\_task\_ids function in internvl/model/internvl\_chat/`modeling\_internvl\_chat.py`, lines 933-979, instead chooses one minimum-error eligible task for each sample, counts those winners over the batch, and returns one or two task IDs based on batch votes and a 0.3 support threshold. The returned IDs are applied through one shared expert mask. The second task requires multiple distinct per-sample winners, so a single sample can never activate two task experts through this function.\par
Minimal input, thresholds (1,1) | Manuscript per-sample rule | Pinned release trace\par
One sample with errors (0.1,0.2) | Both experts 1 and 2 are eligible and selected. | Only expert 1 is returned.\par
Sample x errors (2,2), evaluated alone | No candidate; use pretrained backbone. | No candidate; backbone fallback.\par
Same x batched with a sample having errors (0.1,0.2) | x still has no eligible expert. | Batch votes return expert 1, whose mask also applies to x.\par
These are reviewer-created traces of the visible control flow, not author-code execution. Default evaluation batch sizes are larger than one: the VQA evaluator uses 8, and captioning/grounding use 10; the separate run\_eval.sh specifies 8 or 16. Therefore batch composition is a material part of the released evaluation behavior. A dataset-homogeneous evaluation batch and independent-example deployment need not route the same sample identically.\par
This is a confirmed implementation-to-manuscript discrepancy. Its effect on the original tables remains unverified without the original commit, batch construction, thresholds, routing outputs and experiment logs. It does not establish false scores, test-answer leakage, or a reversed comparison. The inspected evaluation calls supply questions and images to generation, with annotations retained for scoring; no reference-answer leakage was found in those paths.\par
Other code details matter for correspondence: the allocator operates on individual attention linear-module scores, sums these scores by modality, and applies the minimum vision share. The manuscript describes an entire-module gradient L2 norm and no such floor. Exact proxy aggregation and original allocation settings should therefore be documented before claiming that the release reproduces the printed curriculum exactly.\par
The source’s p. 6 statement that any rejecting autoencoder triggers fallback is best read as a local quantifier error: Eq. (10), p. 5 and the per-sample candidate test in the code support an empty-candidate fallback. The batch-mask issue above is separate. A disabled-adapter flag in one class was also examined; the traced main routing paths explicitly clear expert masks before feature extraction, so that flag alone does not establish a defect in those paths.\par
\section*{7. Empirical results: reconstructed gains and limits}
Metric | D-MoLE | Best displayed baseline (O-LoRA) | Absolute gain\par
AVG | 73.87 | 58.79 | 15.08 aggregate-score points\par
Last | 82.18 | 62.04 | 20.14 aggregate-score points\par
BWT | -1.49 | -21.31 | 19.82 aggregate-score points\par
Table 2’s row averages reconstruct within rounding. The quoted 15.08\%, 20.14\% and 19.82\% are arithmetic differences, not relative percentage improvements. AVG and Last mix CIDEr scores with percentage-valued VQA/grounding metrics, so the honest unit is aggregate-score points. Relative changes of the first two displayed aggregates would be about 25.65\% and 32.46\%, although those mixed-unit aggregates still require careful interpretation. This unit clarification does not erase the displayed score advantage.\par
The evaluation definitions are explicit: AVG\_i averages A\_t,i across all N checkpoints, including stages before task i is trained; Last\_i is A\_N,i; BWT\_i averages A\_t,i-A\_i,i over all later checkpoints. They are not the same as an average over only seen tasks or a final-checkpoint-only backward-transfer metric. The paper deliberately evaluates unseen tasks because it starts from a pretrained model. This is a disclosed protocol choice rather than hidden denominator manipulation.\par
The nine datasets span captioning, VQA and grounding, and use common templates within output formats. This reduces task-specific template cues, but does not remove dataset-domain information; Appendix F itself describes dataset boundaries as relatively distinct. A t-SNE plot of reconstructions from each task’s own autoencoder does not by itself establish cross-autoencoder routing accuracy or calibration. Router confusion, unknown-task rejection and batch-stratified results would make that interpretation testable.\par
The v4 ablation removes both dynamic placement and task-specific routing, so its drop cannot isolate the causal contribution of allocation alone. General benchmarks in Table 3 show reduced forgetting relative to the reported adaptation baselines, while MME and MMMU remain below the original zero-shot model. The source supports a comparative retention improvement, not universal preservation of pretrained ability.\par
The threshold sensitivity table evaluates a fixed final checkpoint and gives aggregate Last scores from 80.34 to 82.18 across the shown multipliers. It supports limited aggregate stability under those settings; individual tasks vary more. Exact thresholds, original calibration data and raw predictions remain unavailable. The author’s threshold-calibration reply (repository issue 2) recommends setting thresholds from training reconstruction-loss distributions, for example above a high percentile; this guidance is not the original run’s threshold file.\par
The reported preprocessing and routing overheads were inspected rather than assumed absent. Table 9 sums to 44,772 training seconds and 650 proxy-preprocessing seconds (1.4518\%); including both gives 12.6172 hours. Table 10’s 0.0552 s routing versus 0.5225 s inference gives 10.5646\%, matching its 10.56\%. In Table 8, listed task-time sums differ from reported totals by 0.08 h for Seq-FT, 0.07 h for EWC-LoRA and 0.09 h for Dense MoLE, exceeding a 0.05 h independent-rounding allowance. D-MoLE’s 44,772 s is 12.4367 h rather than the displayed 12.40 h. Different timing boundaries or runs could explain these differences; logs and inclusion of autoencoder training are needed. An IconQA epoch discrepancy between Appendices A and H also remains. These issues do not establish a reversal of the broad training-time ordering.\par
No repeated-run significance test or raw checkpoint-by-task score matrix was reproduced. The word significantly should not be interpreted as an independently verified statistical result. The public issue alleging different reproduced values is an unverified user report; it is not treated as evidence that the paper’s scores are wrong.\par
\section*{Supporting checks}
The publisher record is PMLR 267:19011-19033. The available arXiv record is 2506.11672v1, submitted 2025-06-13, with a separately preserved 24-page PDF (SHA-256 `d26c6fe9612c637d6cad7e40ed9dbff50074c48cb5034231d72767edb70a02d2`). The full files differ, but Appendix K is identical after whitespace removal: 3,287 characters, SHA-256 `be4e986cb117d67a31c889d879f97f8dda92c52cce4567268dbdcf61b5d87c25`. No verified correction excluding the counterexample was found in the inspected primary sources.\par
The repository’s current pinned commit is dated 2026-01-11; its most recent changes concern download links, with the three checked core files (router, allocator and scorer) unchanged since the 2025-06-16 initial release. Forty-six relevant source text files were preserved and their Git blob hashes verified. The audit also preserves repository metadata, public issue context and official base-model configuration evidence. A contemporaneous public release is not proof that it produced Table 2. The arXiv comparison targeted Appendix K; no complete table-by-table version audit is claimed.\par
Executed work was limited to reviewer-created analytic arithmetic, table sums, source-figure measurements, source acquisition and PDF processing. No author training/evaluation/scoring code was run or imported, and no neural models were run.\par
The scientific decision remains pending; human verification and author adoption are not asserted.\par
\section*{9. Conclusion and primary references}
Issue | Supported status | Required action\par
Theorem 3.1 / K.1 | Major: a complete admissible counterexample refutes the formal claim. | Remove/weaken the theorem or supply substantive new assumptions and a valid proof.\par
Published parameter-count equality | Major: the stated equal-budget rationale fails; original checkpoint counts and score effects remain Needs Evidence. | Provide per-task architecture/tensor counts and reconcile attention-only versus released baseline targets.\par
Per-sample versus batch routing | Confirmed release discrepancy; effect on original numerical claims unverified. | Supply original implementation/protocol and test sample-level versus batch-vote routing.\par
Headline score gains | Displayed arithmetic reconstructs; percent wording has the wrong unit. | Report aggregate-score point differences and explicit metric definitions.\par
[P1] Ge et al. (2025), Dynamic Mixture of Curriculum LoRA Experts for Continual Multimodal Instruction Tuning. https://proceedings.mlr.press/v267/ge25d.html . Controlling source: papers/0085/0085.pdf.\par
[P2] Ge et al., arXiv:2506.11672v1. https://arxiv.org/abs/2506.11672 . Relevant-section comparison and separate PDF hash are recorded above.\par
[P3] Official implementation: https://github.com/gcd19/D-MoLE/tree/525fe2b788da85ebe3f5f102c4a77b3d37348c18 . Source manifest: verification/0085/`agent\_code\_selected\_tree.json`.\par
[P4] Official InternVL2-2B configuration at revision ad4b36dc0ac1bd114570c5759dac3a04cfa573d5. Pinned InternVL2-2B configuration (config.json) (https://huggingface.co/OpenGVLab/InternVL2-2B/blob/ad4b36dc0ac1bd114570c5759dac3a04cfa573d5/config.json) .\par
The norm theorem and the published matched-budget rationale each have a demonstrated core dependency. The routing discrepancy has only a verified release-level scope. Historical parameter counts, corrected matched-budget results and the observed score ordering require the original experiment artifacts for independent verification.\par
\end{document}
